\documentclass[11pt,a4paper]{article}
\usepackage[margin=25mm]{geometry}
\usepackage{amsmath,amssymb,amsthm,booktabs}
\newtheorem{lemma}{Lemma}\newtheorem{theorem}{Theorem}
\title{Affine subgroups and the remaining counts with an entry of order eight}
\author{Human Compression project}\date{}
\begin{document}\maketitle
Let $G=O'N$ and $g=|G|=460815505920$. Write $x_T(p,q)$ for the
number of ordered pairs $(r,s)$ in $T$ with $|r|=p$, $|s|=q$ and
$|rs|=2$, and $e_T(p,q)$ for the number generating $T$. We prove
the two remaining order-eight entries of the fixed Connor--Leemans table:
\[
 e_G(4,8)/(2g)=284,\qquad e_G(8,8)/(2g)=5424.                 \tag{1}
\]
There is no quotient by reversal. The complete subgroup roster,
the acquired character evaluations, and the named classical groups
are explicit borrowed premises. Row numbers always refer to that roster.
The small permutation computations below prove finite statements by
exhaustive multiplication in groups of order at most 25920. Identification
of each model is proved separately; agreement of order spectra is only
a check.

\section*{The involution centralizers}
Let $\mathcal A$ be the class of maximal involution centralizers, row 7.
The accepted order-eight centralizer lemma says that for every $|s|=8$
the three involutions in $C_G(s)$ form a four-group $V_s$ of row 578.
Its centralizer has order 256 and its normalizer has order 1536.
The accepted literal reader gives
\[
 x_{C_G(V_s)}(4,8)=3072,\qquad x_{C_G(V_s)}(8,8)=672.
\]
Consequently the number $U_{\mathcal A}$ of pairs covered by this
maximal family is
\[
 U_{\mathcal A}(p,8)=[G:A]x_A(p,8)
              -2[G:N_G(V_s)]x_{C_G(V_s)}(p,8).              \tag{2}
\]
A pair has zero, one or three common commuting involutions, which
proves (2) directly. This is also the rank-two case of the accepted
elementary abelian subgroup identity. Substitution gives
\[
 U_{\mathcal A}(4,8)/g=53/2,\qquad
 U_{\mathcal A}(8,8)/g=983/16.                              \tag{3}
\]

\section*{The affine group with an alternating factor}
Put $B$ for row 11, of order 25920. The accepted structural argument
identifies its normal subgroup of row 17 as
$B_0=(V_9:C_4)\times N$, where $V_9\cong C_3^2$, $N\cong A_6$,
and $C_4$ acts faithfully on $V_9$. Its subgroup $H$ of row 32 has
order 2880, $H/N\cong Q_8$, and
$H\cap B_0=N\times C_4$. These are structural inputs already proved
in the order-ten and order-twelve components.

\begin{lemma}\label{B}
The group $B$ is the affine pullback
\[
 V_9\rtimes\bigl(M_{10}\mathbin{\times}_{C_2}Q_8\bigr),     \tag{4}
\]
where $Q_8$ acts faithfully on $V_9$ and its specified cyclic subgroup
of order four is the kernel of the map to $C_2$. The two index-two
subgroups in rows 18 and 19 are both isomorphic to
\[
 D=V_9\rtimes\bigl(M_{10}\mathbin{\times}_{C_2}C_4\bigr),  \tag{5}
\]
with the faithful action of $C_4$ on $V_9$.
\end{lemma}
\begin{proof}
Both $V_9=O_3(B_0)$ and the last derived subgroup $N$ of $B_0$ are
normal in $B$. Their intersection is trivial and they commute.
The given intersection with $H$ implies $H\cap V_9=1$, so orders
give $B=V_9\rtimes H$. The map from $H$ into
$\operatorname{Aut}(N)\times Q_8$ given by conjugation and quotient
is injective: its kernel is contained in $N$ and centralizes $N$.
Its image is the full pullback of the induced map
$\phi:Q_8\longrightarrow\operatorname{Out}(A_6)$.
The specified cyclic subgroup is in $\ker\phi$.

The classical description of $\operatorname{Aut}(A_6)$ has three
index-two groups, $S_6$, $\operatorname{PGL}_2(9)$ and $M_{10}$.
We use this standard classification explicitly, with the primary
ATLAS names and the retained classical automorphism theorem.
If $\phi$ is trivial, or its image gives $S_6$, choose an $A_5$
subgroup of $N$ in a class fixed by that image. Its normalizer in
$H$ has order 480 and maps onto $Q_8$: in the trivial case it is
$A_5\times Q_8$; in the $S_6$ case it is the pullback of $S_5$.
Thus its normalizer in $B$ is $V_9\rtimes N_H(A_5)$, of order 4320,
and still maps onto $Q_8$. No maximal subgroup in the complete
row 11 list can contain it. The three maximal subgroups of order
12960 have proper image in $Q_8$, and the other maximal orders
2880, 2592, 720 and 576 are not divisible by 4320.

If the image gives $\operatorname{PGL}_2(9)$, take an outer element
of order ten in its nonsplit torus. A lift paired with a quaternion
element outside $\ker\phi$ has order twenty and is outside $B_0$.
But the acquired character table gives exactly 2592 elements of
order twenty in $B$. All are already in $B_0$: its affine factor
has eighteen elements of order four and $A_6$ has 144 elements
of order five. This excludes the second possibility. Therefore
the image is $M_{10}$ and its kernel is precisely the specified
$C_4$.

The action of $Q_8$ on $V_9$ is faithful, since any nontrivial normal
subgroup of $Q_8$ contains its central involution, which acts
nontrivially in the specified $C_4$. The faithful two-dimensional
quaternion module over $\mathbb F_3$ is unique: the four quaternion
matrix monomials span $M_2(\mathbb F_3)$. The literal reader checks
this and that $\operatorname{GL}_2(3)$ realizes all 24 automorphisms
of the quaternion subgroup. In particular the two remaining cyclic
subgroups of order four can be exchanged while preserving the
specified one. This proves (4) and (5), including the assignment
of the character to the chosen cyclic subgroup.
\end{proof}

\begin{lemma}\label{Bcount}
The union of the row 11 family outside $\mathcal A$ contributes
$32g$ pairs of type $(4,8)$ and $48g$ pairs of type $(8,8)$.
\end{lemma}
\begin{proof}
If the images of $r,s$ generated $Q_8$, they would have to be two
elements from distinct cyclic subgroups of order four; their product
would then have order four, contrary to $|rs|=2$. Hence they lie in
one cyclic subgroup. The specified preimage $B_0$ has no element of
order eight. In particular $s$ cannot map to the central subgroup of
order two, whose preimage is inside $B_0$. Every pair therefore lies
in exactly one of the two subgroups (5).

For the explicit group $D$, the complete finite calculation gives
\[
\begin{array}{c|rrr}
p&x_D(p,8)&\text{common internal involution}&e_D(p,8)\\\hline
4&648000&233280&414720\\
8&839160&217080&622080.
\end{array}                                                \tag{6}
\]
Every pair not generating $D$ has a common internal involution.
The finite proof partitions the four classes of elements of order
eight, each of size 810, then lists the orbits of all partners under
the full centralizer of order sixteen. For each representative it
closes the two generators by literal multiplication. The generated
orders are respectively $32,288,1440,12960$ and
$8,16,144,1440,12960$; all proper cases have the stated involution.
Conversely a generating pair of $D$ cannot be in a member of
$\mathcal A$, since $12960$ does not divide $161280$.
The two generated subgroup classes are exactly rows 18 and 19,
each with normalizer order 25920. Counting by the subgroup actually
generated therefore gives $2g\,e_D/25920$, which is $32g$ or $48g$.
Distinct generated subgroups give disjoint sets of pairs.
\end{proof}

\section*{The affine group with an extraspecial factor}
Put $C$ for row 12. Its borrowed primary structure is $V\rtimes K$,
where $V\cong C_3^4$, $|K|=320$, and $K$ has a normal extraspecial
subgroup $E$ of order 32 with quotient $D_{10}$. The accepted
extraspecial argument identifies the minus type: the five nonzero
singular vectors of $E/Z(E)$ are cyclically permuted by an element
of order five.

\begin{lemma}\label{Cmodel}
The faithful four-dimensional affine model supplied below is
isomorphic to $C$. Its linear complement is the full preimage of
a dihedral subgroup of order ten in the permutation action on the
five singular vectors of $E/Z(E)$.
\end{lemma}
\begin{proof}
The central involution of $E$ cannot centralize $V$: otherwise its
centralizer in $G$, of order 161280, would contain a subgroup of
order 81. A nontrivial normal subgroup of an extraspecial group
meets its center, so $E$ acts faithfully on $V$.

Here is an elementary model and uniqueness argument for this module.
Over $\mathbb F_3$ put
\[
 a=\begin{pmatrix}1&0\\0&-1\end{pmatrix},\quad
 b=\begin{pmatrix}0&1\\1&0\end{pmatrix},\quad
 i=\begin{pmatrix}0&1\\-1&0\end{pmatrix},\quad
 j=\begin{pmatrix}1&1\\1&-1\end{pmatrix}.
\]
The matrices $a\otimes I,b\otimes I,ab\otimes i,ab\otimes j$
are four pairwise anticommuting involutions. Their product is a
fifth involution with the same property. They generate $E$ with
central involution $-I$. Their sixteen square-free monomials form
a basis of $M_4(\mathbb F_3)$, as the reader checks by exact row
reduction. The quotient of $\mathbb F_3[E]$ in which the central
involution acts as $-1$ has dimension sixteen, and is thus precisely
this matrix algebra. Every nonzero module for this algebra has
dimension at least four. By Maschke's theorem the nonzero minus
eigenspace in the actual four-dimensional module is therefore
the whole module. This proves uniqueness and shows that the center
acts as $-I$ on all of $V$.

The action of $K/E$ on $E/Z(E)$ is faithful. Any nontrivial kernel
contains the rotation subgroup of order five. Its preimage has an
element of order five by Cauchy's theorem; that element acts on $E$
by an inner automorphism and hence trivially, since the inner
automorphism group is a 2-group. This would give an element of order
twenty in $K$, absent from the acquired spectrum of $C$.
It follows that $C_K(E)=Z(E)$: after its image in $K/E$ vanishes,
it lies in the center of $E$. Any kernel of the action on $V$
centralizes $E$ and meets $E$ trivially, so is trivial.

The normalizer of the explicit $E$ in $\operatorname{GL}_4(3)$
acts on the five singular vectors. Its kernel is $E$: an automorphism
trivial on $E/Z(E)$ is inner, and the matrix centralizer is the
scalar group already in $E$. Explicit intertwining matrices give
a five-cycle and a transposition on these five vectors. Hence the
normalizer induces $S_5$. All dihedral subgroups of order ten in
$S_5$ are conjugate, and their full preimages have order 320.
The displayed matrices for a five-cycle and a reflection generate
one such preimage together with $E$. This identifies the actual
linear complement and hence the affine group, without inferring
isomorphism merely from an order spectrum.
\end{proof}

\begin{lemma}\label{Ccount}
Apart from pairs generating a subgroup in rows 250, 252, 332 or 333,
the row 12 family outside $\mathcal A$ contributes $8g$ pairs of
type $(4,8)$ and $20g$ pairs of type $(8,8)$. Each such pair has a
unique row 12 maximal overgroup and no other maximal overgroup.
\end{lemma}
\begin{proof}
An element of order eight has nontrivial reflection image in
$K/E\cong D_{10}$, since $E$ has exponent four and $V$ has odd
order. Two elements of 2-power order with involutory product have
image in the same reflection subgroup. Thus a pair lies in exactly
one of the five preimages $D_C$ of order 5184, the subgroups of row 26.
In the explicit affine group the complete generated-order partitions
are
\[
\begin{array}{c|rrrrr}
p=4&16&32&144&1296&2592\\
\text{pairs}&5184&5184&20736&20736&20736\\\hline
p=8&8&16&72&648&1296\\
\text{pairs}&1296&2592&20736&82944&20736.
\end{array}                                                \tag{7}
\]
The two classes of order-eight elements each have size 648 and
centralizer order eight. Complete partner orbits and generated
closures give (7). Precisely the cases of 2-power generated order
have a common internal involution. In the cases of orders divisible
by 81, the generated subgroup contains $V$: its intersection with
$V$ is its unique Sylow 3-subgroup. Since the sole order-81 class
has normalizer 25920, every row 12 overgroup is $N_G(V)$ and is
unique. The complete roster ancestors of these possible generated
classes give no other maximal family. These pairs cannot be in
$\mathcal A$, whose members have 3-part nine.

The remaining orders 144 and 72, with an element of order eight,
are exactly rows 250, 252, 332 and 333 among the descendants of row 12.
They will be counted with the projective and Mathieu families.
The full-$V$ counts in (7), multiplied by five and by $g/25920$,
give $8g$ and $20g$ as stated.
\end{proof}

\section*{The projective and Mathieu families}
Write $M=L_3(7):2$, in either of rows 2 and 3, and $J=M_{11}$,
in either of rows 23 and 24. The accepted block and orthogonal
centralizers in $M$ are $H$ of order 1344 and $K_0$ of order 672.
There are 2793 inner involutions and 5586 outer involutions.

\begin{lemma}\label{Mcover}
The pairs in $M$ covered by $\mathcal A$ number
\[
 2793\bigl(x_H(p,8)-T_H(p,8)\bigr)+5586x_{K_0}(p,8),        \tag{8}
\]
where $T_H(4,8)=0$, $T_H(8,8)=1008$. Thus, summing the two maximal
classes, the pairs outside $\mathcal A$ number $56g$ and $272g$.
\end{lemma}
\begin{proof}
For any $s\in M$ of order eight, its fourth power is an inner
involution. Thus $s$ lies in its block centralizer $H$.
The three order-eight classes of the complete block model have
sizes 42, 42 and 84. Their centralizers have orders 32, 32 and 16,
but each contains exactly three involutions, one inner and two outer.
By the global centralizer lemma these are all the involutions of
$C_G(s)$. Hence every common commuting involution of a pair in $M$
already belongs to $M$.

A four-group centralized by $s$ has exactly one inner involution,
by the same literal class calculation. Counting involutions gives
the positive terms $2793x_H+5586x_{K_0}$ in (8). For a fixed inner
involution $t$, the four-groups containing it are described by the
outer involutions $b\in C_M(t)$, with $b$ and $bt$ describing the
same group. The usual coefficient minus two is therefore exactly
minus the sum over these $b$ of pairs centralizing $b$. This is
$-T_H$, and the retained model gives the displayed values by testing
every partner for all three order-eight classes. Substitution of
\[
 (x_H(4,8),x_H(8,8))=(672,504),\qquad
 (x_{K_0}(4,8),x_{K_0}(8,8))=(2688,3192)
\]
gives (8). Subtracting it from the acquired character count of $M$
and multiplying by its two class lengths gives $56g$ and $272g$.
\end{proof}

The accepted complete $M_{11}$ model gives the following refinement
outside its internal involution centralizers:
\[
\begin{array}{c|rrrr}
 &\text{generated order }72&144&720&7920\\\hline
(4,8)&0&15840&31680&0\\
(8,8)&15840&0&47520&31680.
\end{array}                                                \tag{9}
\]
These generated classes are exactly 332/333, 250/252, 102/103 and
23/24, respectively. None is a descendant of row 7, so internal
and global involution coverage agree for these pairs. The complete
Mathieu families outside $\mathcal A$ therefore contribute $12g$
and $24g$ before correcting the overlap with the projective families.

\begin{lemma}\label{overlap}
The projective and Mathieu unions outside $\mathcal A$ intersect
in exactly $4g$ pairs for each of the two types. No pair is counted
twice within either family.
\end{lemma}
\begin{proof}
The complete common descendant lists, restricted to subgroups
containing an element of order eight and excluding row 7 descendants,
are precisely rows 250, 252, 332 and 333. A row 250 or 252 subgroup
has exactly one actual projective and one actual Mathieu overgroup:
the direct class-length incidences with its listed maximal parents
are one. Its other parent belongs only to the row 12 family.
A row 332 or 333 subgroup has a unique actual parent in row 252
or 250, respectively, and its other parent also belongs only to
row 12. Hence the same uniqueness holds for order 72. Formula (9)
shows that precisely $15840$ pairs per $M_{11}$ have one of these
generated subgroups. The two class lengths give
$2g\cdot15840/7920=4g$.

For completeness, outside these four rows and the affine families,
the remaining projective generated classes are
\[
 2,3,4,5,9,10,13,14,45,46,53,54,91,92,93,94,95,96,154,155,172,174.
\]
The accepted parabolic alignment lemma gives a unique actual
maximal overgroup for rows 91, 92, 154 and 155. Every other
nonmaximal row has the previously proved unique immediate-parent
chain to row 2 or 3. That lemma depends on the subgroup class and
the presence of an element of order eight, not on an order-seven
entry in the generating pair. It applies unchanged here.
The possible Mathieu generated classes in (9) have the uniqueness
already stated, or generate the maximal itself. No arbitrary
product of cover multiplicities is being used.
\end{proof}

\section*{The final subtraction}
The row 21 family $4^3.L_3(2)$ contributes no pairs outside
$\mathcal A$, by the accepted complete literal calculation.
The character evaluations give zero in every other maximal family
except those already treated and the two classes of $L_2(31)$.
In $L=L_2(31)$ every proper subgroup containing an element of order
eight lies in its dihedral torus normalizer of order 32. For type
$(4,8)$ no such pair has product of order two. For type $(8,8)$
the factors must be in the cyclic rotation group, and there are
four pairs in its unique cyclic subgroup of order eight. There
are $|L|/32$ such normalizers. The acquired raw counts are
$4|L|$ and $33|L|/8$ respectively, so
\[
                         e_L(4,8)=e_L(8,8)=4|L|.
\]
All proper pairs lie in a 2-group and therefore have a common
central involution. The two generated maximal classes contribute
$8g$ in each case.

Combining (3), (6)--(9) and the overlap gives
\[
\begin{array}{c|rr}
\text{contribution divided by }g&(4,8)&(8,8)\\\hline
\mathcal A&53/2&983/16\\
\text{row 11, outside }\mathcal A&32&48\\
\text{row 12 with full }V&8&20\\
\text{projective families, outside }\mathcal A&56&272\\
\text{Mathieu families, outside }\mathcal A&12&24\\
\text{projective--Mathieu intersection}&-4&-4\\
\text{generated }L_2(31)&8&8\\\hline
\text{all nongenerating pairs}&277/2&6871/16.
\end{array}                                                \tag{10}
\]
Every subtraction has been attached to an actual common subgroup
or to the subgroup generated by the pair. The raw character totals
are $1413g/2$ and $180439g/16$. Subtracting (10), and dividing by
$2g=|\operatorname{Aut}(G)|$, proves (1).

\section*{Attribution and finite verification}
The numerical targets are the earlier published entries. The source
roster, primary group structures and acquired character tables are
borrowed, not newly classified here. The accepted common-involution
formula and parabolic alignment are reused. The new work is the
explicit affine model identification, generated-subgroup partitions,
and the common subtraction in (10). The literal readers verify
finite permutation and matrix identities; they are not Lean proofs.
This component alone makes no assertion about chirality or completion
of the entire original.
\end{document}
